%%%%%%%%%%%%%%%%%%%%%%%%%%%%%%%%%%%%%%%%%
% کارنامۀ عمومی (نسخۀ فارسی)
% معادل فارسیِ ../Kahbod-Aeini_CV.tex
% کامپایل: xelatex Kahbod-Aeini_CV_FA.tex
%%%%%%%%%%%%%%%%%%%%%%%%%%%%%%%%%%%%%%%%%

\documentclass[11pt]{article}
\usepackage{fa-cv}

\begin{document}

\faHeader{کهبد آیینی}{%
دانشکدۀ مهندسی برق و کامپیوتر، دانشگاه واترلو، واترلو، انتاریو، کانادا\\[0.3em]
\lr{\href{mailto:kaeini@uwaterloo.ca}{kaeini@uwaterloo.ca}} |
\lr{\href{https://www.linkedin.com/in/kahbod-aeini-833b71190/}{LinkedIn}} |
\lr{\href{https://github.com/kahbodaeini}{GitHub}} |
\lr{\href{https://kahbod.org}{kahbod.org}} |
\lr{\href{https://scholar.google.ca/citations?user=3Fd5R9YAAAAJ&hl=en}{Google Scholar}}%
}

%----------------------------------------------------------------------------------------
%	تحصیلات
%----------------------------------------------------------------------------------------
\begin{faSection}{تحصیلات}

\faEntry{دانشگاه واترلو}{ژانویۀ ۲۰۲۶ تاکنون}
دکتری مهندسی برق و کامپیوتر\\
استاد راهنما: پروفسور \lr{\href{https://uwaterloo.ca/distributed-algorithms-systems-lab/}{Wojciech Golab}}

\vspace{0.5em}
\faEntry{دانشگاه کلگری}{۲۰۲۳ تا ۲۰۲۵}
کارشناسی ارشد علوم کامپیوتر\\
استاد راهنما: پروفسور \lr{\href{https://www.woelfellab.org/about.html}{Philipp Woelfel}}\\
پایان‌نامه: \href{https://ucalgary.scholaris.ca/bitstreams/5ae5cf8f-5009-4db3-b91e-d3dbb679b5bf/download}{«قفل بی‌انتظارِ تصادفی با فضای بهینه»}

\vspace{0.5em}
\faEntry{دانشگاه صنعتی شریف (آریامهر)}{۲۰۱۹ تا ۲۰۲۳}
کارشناسی مهندسی کامپیوتر\\
پایان‌نامه: \href{https://github.com/kahbodaeini/Fraud-Detection}{«تشخیص تقلب از روی دنباله‌های تراکنش‌های بانکی با الگوریتم‌های گوناگون یادگیری ماشین»}

\end{faSection}

%----------------------------------------------------------------------------------------
%	علایق پژوهشی
%----------------------------------------------------------------------------------------
\begin{faSection}{علایق پژوهشی}
سیستم‌های توزیع‌شده، الگوریتم‌های همروند و تصادفی، حافظۀ پایدار و اشیای بازیابی‌پذیر،
سیستم‌های حافظۀ مشترک، برنامه‌نویسی چندهسته‌ای، ساختمان‌های دادۀ پیشرفته، سیستم‌های پایگاه داده
\end{faSection}

%----------------------------------------------------------------------------------------
%	انتشارات
%----------------------------------------------------------------------------------------
\begin{faSection}{انتشارات}

\textbf{کهبد آیینی}، \lr{Dante Bencivenga}، \lr{George Giakkoupis}، \lr{Philipp Woelfel}\\
\href{https://dl.acm.org/doi/10.1145/3796701.3815952}{\textit{«قفل‌های بی‌انتظارِ تصادفیِ ساده و کارآمد»}}\\
\href{https://www.podc.org/podc2026/}{چهل‌وپنجمین سمپوزیوم \lr{ACM} در مبانی محاسبات توزیع‌شده (\lr{PODC})}، ۲۰۲۶. \textbf{جایزۀ بهترین مقاله.}

\vspace{0.5em}
\textbf{کهبد آیینی}، \lr{Wojciech Golab}\\
\href{https://dl.acm.org/doi/10.1145/3820355.3820388}{\textit{«تحلیل خطی‌پذیری در سیستم‌های توزیع‌شدۀ نسبیتی»}}\\
هشتمین کارگاه \lr{ApPLIED}، هم‌زمان با \lr{PODC 2026}، ۲۰۲۶.

\vspace{0.5em}
\textbf{کهبد آیینی}، \lr{Brianna Wang}، \lr{Wojciech Golab}، \lr{Neeraj Mittal}\\
\textit{«ساخت جهانیِ سریع‌ترِ اشیای خواندن-تغییر-نوشتنِ آشکارپذیر برای حافظۀ پایدار»}\\
سمپوزیوم \lr{ACM SIGPLAN} در اصول و عمل برنامه‌نویسی موازی (\lr{PPoPP})، ۲۰۲۷. \textbf{در دست داوری.}

\vspace{0.5em}
\textbf{کهبد آیینی}\\
\textit{«قفل بی‌انتظارِ تصادفی با فضای بهینه»}\\
\href{https://ucalgary.scholaris.ca/bitstreams/5ae5cf8f-5009-4db3-b91e-d3dbb679b5bf/download}{پایان‌نامۀ کارشناسی ارشد}، دانشگاه کلگری، ۲۰۲۵.

\end{faSection}

%----------------------------------------------------------------------------------------
%	سوابق کاری
%----------------------------------------------------------------------------------------
\begin{faSection}{سوابق کاری}

\faEntry{مهندس نرم‌افزار، \lr{ArcTrade} (پاره‌وقت)}{دسامبر ۲۰۲۵ تاکنون}
\begin{faItems}
  \item توسعۀ سرویس‌های سمت سرور و واسط‌های برنامه‌نویسی (\lr{API}) برای سامانۀ عامل خرده‌فروشی انرژیِ \lr{ArcTrade}.
  \item پیاده‌سازی تبادل الکترونیکی دادۀ امن و خودکار (\lr{EDI}) با شرکای تجاری بیرونی، برای مبادلۀ رمزنگاری‌شده و تأییدپذیر تراکنش‌های کسب‌وکار.
  \item مشارکت در بخش‌های پشتیبانِ معاملات انرژی و مدیریت ریسک، قیمت‌گذاری، صورت‌حساب و تسویه در بازارهای برق آمریکای شمالی.
  \item همکاری با تیم‌های میان‌رشته‌ای برای طراحی راهکارهای نرم‌افزاری مقیاس‌پذیر و کارآمد.
\end{faItems}

\faEntry{مهندس داده، \lr{AI Lida}}{فوریه تا سپتامبر ۲۰۲۲}
\begin{faItems}
  \item توسعه، نگهداری و بهینه‌سازی خطوط لولۀ داده برای فرایندهای \lr{ETL} با \lr{Python} و \lr{PostgreSQL}.
  \item همکاری با واحد امنیت برای گردآوری نیازمندی‌ها و ارائۀ راهکارهای دادهایِ هم‌راستا با اهداف کسب‌وکار.
  \item مصاحبه، استخدام، آموزش و سرپرستی تیمی پنج‌نفره از کارآموزان مهندسی داده، همراه با راهنمایی و بازخورد عملکرد.
\end{faItems}

\end{faSection}

%----------------------------------------------------------------------------------------
%	سوابق پژوهشی
%----------------------------------------------------------------------------------------
\begin{faSection}{سوابق پژوهشی}

\faEntry{دانشجوی دکتری و دستیار پژوهشی، الگوریتم‌های توزیع‌شده --- دانشگاه واترلو}{ژانویۀ ۲۰۲۶ تاکنون}
استاد راهنما: پروفسور \lr{Wojciech Golab}\\
پژوهش در حوزۀ حافظۀ پایدار، سیستم‌های توزیع‌شدۀ مبتنی بر حافظۀ مشترک و اشیای بازیابی‌پذیر، با تمرکز بر طراحی و تحلیل الگوریتم‌های کارآمد برای همگام‌سازی، اجماع و تحمل خطا در سیستم‌های توزیع‌شده.

\vspace{0.5em}
\faEntry{دانشجوی کارشناسی ارشد و دستیار پژوهشی، الگوریتم‌های توزیع‌شده --- دانشگاه کلگری}{ژوئن ۲۰۲۴ تا ژوئن ۲۰۲۵}
استاد راهنما: پروفسور \lr{Philipp Woelfel}\\
طراحی و تحلیل الگوریتم‌های قفلِ بی‌انتظارِ تصادفی با موازنۀ بهینۀ فضا و زمان در سیستم‌های حافظۀ مشترک؛ از جمله نخستین الگوریتم با فضای بهینه و گونه‌های مکمل با زمان بهینه.

\vspace{0.5em}
\faEntry{دستیار پژوهشی زیست‌اطلاعاتی (داوطلبانه، از راه دور) --- دانشگاه جانز هاپکینز}{فوریه تا اوت ۲۰۲۳}
استاد راهنما: پروفسور \lr{Mehdi Pirooznia}\\
انتخاب ویژگی و توسعۀ مدل‌های یادگیری ماشین برای پیش‌بینی ترکیب سلول‌های ایمنی در سرطان ریه؛ شناسایی انواع سلول‌های ایمنی، تحلیل ژن‌های با بیان متفاوت، تحلیل غنی‌سازی عملکردی و دسته‌بندی زیرگونه‌های سرطان با روش‌های یادگیری ماشین.

\vspace{0.5em}
\faEntry{دستیار پژوهشی یادگیری عمیق --- پژوهشگاه دانش‌های بنیادی (\lr{IPM})}{اکتبر ۲۰۲۱ تا فوریۀ ۲۰۲۲}
اساتید راهنما: دکتر فاطمه بهاری‌فرد و دکتر وحید متقد\\
توسعه و ارزیابی الگوریتم‌های خوشه‌بندی چهره با شبکه‌های کانولوشنی گرافی (\lr{GCN}) برای پیش‌بینی پیوند و دسته‌بندی زوجی.

\end{faSection}

%----------------------------------------------------------------------------------------
%	جوایز و افتخارات
%----------------------------------------------------------------------------------------
\begin{faSection}{جوایز و افتخارات}

\faEntry{جایزۀ بهترین مقاله --- چهل‌وپنجمین سمپوزیوم \lr{ACM} در مبانی محاسبات توزیع‌شده (\lr{PODC})}{۲۰۲۶}
برای مقالۀ «قفل‌های بی‌انتظارِ تصادفیِ ساده و کارآمد».

\vspace{0.4em}
\faEntry{جایزۀ پژوهشی دانشکده --- دانشکدۀ علوم کامپیوتر، دانشگاه کلگری}{۲۰۲۳}
جایزۀ دستیاری پژوهشی که به بهترین دانشجویان هر دانشکده تعلق می‌گیرد تا به‌جای وظایف آموزشی بر پژوهش خود تمرکز کنند.

\vspace{0.4em}
\faEntry{میان صد نفر برتر کشور}{۲۰۱۹}
رتبۀ ۹۴ کشوری (۳۱ منطقه) در کنکور سراسری میان حدود ۱۶۰٬۰۰۰ شرکت‌کننده.

\vspace{0.4em}
\faEntry{بورسیۀ بنیاد ملی نخبگان}{۲۰۱۹}
عضویت و بورسیه در بنیاد ملی نخبگان ایران.

\end{faSection}

%----------------------------------------------------------------------------------------
%	سوابق تدریس
%----------------------------------------------------------------------------------------
\begin{faSection}{سوابق تدریس}

\faEntry{دانشگاه کلگری}{۲۰۲۳ تا ۲۰۲۶}
سرپرستی و مدیریت تیم‌های دستیار آموزشی در درس‌های کارشناسی و تحصیلات تکمیلی (در مجموع تا ۴۳۰ دانشجو)؛ از جمله مبانی نظری علوم کامپیوتر (۷ دستیار آموزشی، بیش از ۱۵۰ دانشجو)، الگوریتم‌های تصادفی (تحصیلات تکمیلی)، پارادایم‌های برنامه‌نویسی (\lr{Haskell} و \lr{Prolog}) و برنامه‌نویسی \lr{Python}.

\vspace{0.5em}
\faEntry{دانشگاه صنعتی شریف}{۲۰۲۰ تا ۲۰۲۳}
سرپرست دستیاران آموزشی در چند درس کارشناسی (تیم‌های ۱۰ تا ۲۵ نفره، بیش از ۲۰۰ دانشجو)؛ از جمله ساختمان داده و الگوریتم‌ها، محاسبات عددی و برنامه‌نویسی موبایل.

\end{faSection}

%----------------------------------------------------------------------------------------
%	مهارت‌ها
%----------------------------------------------------------------------------------------
\begin{faSection}{مهارت‌ها}

\textbf{زبان‌های برنامه‌نویسی}: \lr{C\#}، \lr{Python}، \lr{R}، \lr{Haskell}، \lr{Java}، \lr{C/C++}، \lr{Prolog}

\vspace{0.3em}
\textbf{چارچوب‌ها و ابزارها}: \lr{.NET / ASP.NET Core}، \lr{Entity Framework}، \lr{Docker}، \lr{Git}، \lr{Django}، \lr{PyTorch}، \lr{TensorFlow}، \lr{scikit-learn}، \lr{NumPy}، \lr{Pandas}، \lr{OpenCV}، \lr{Android}، \lr{limma}

\vspace{0.3em}
\textbf{پایگاه‌های داده}: \lr{Microsoft SQL Server}، \lr{PostgreSQL}، \lr{MongoDB}، \lr{MySQL} \quad \textbf{سایر}: \lr{Linux}، \lr{NixOS}، \lr{Kali}

\end{faSection}

\end{document}
